\begin{lemma}
\label{lemma-fitting-ideals-and-pd1}
Let $R$ be a ring. Let $M$ be a finitely presented $R$-module. Let $k \geq 0$.
Assume that $\text{Fit}_k(M) = (f)$ for some nonzerodivisor $f \in R$
and $\text{Fit}_{k - 1}(M) = 0$. Then
\begin{enumerate}
\item $M$ has projective dimension $\leq 1$,
\item $M' = \Ker(f : M \to M)$ is the $f$-power torsion submodule of $M$,
\item $M'$ has projective dimension $\leq 1$,
\item $M/M'$ is finite locally free of rank $k$, and
\item $M \cong M/M' \oplus M'$.
\end{enumerate}
\end{lemma}

\begin{proof}
Choose a presentation
$$
R^{\oplus m} \xrightarrow{A} R^{\oplus n} \to M \to 0
$$
for some matrix $A$ with coefficients in $R$.

\medskip\noindent
We first prove the lemma when $R$ is local.
Set $M' = \{x \in M \mid fx = 0\}$ as in the statement. By
Lemma \ref{lemma-principal-fitting-ideal} we can choose
$x_1, \ldots, x_k \in M$ which generate $M/M'$. Then
$x_1, \ldots, x_k$ generate $M_f = (M/M')_f$. Hence,
if there is a relation $\sum a_ix_i = 0$ in $M$,
then we see that $a_1, \ldots, a_k$ map to zero in $R_f$
since otherwise $\text{Fit}_{k - 1}(M) R_f = \text{Fit}_{k - 1}(M_f)$
would be nonzero.
Since $f$ is a nonzerodivisor, we conclude $a_1 = \ldots = a_k = 0$.
Thus $M \cong R^{\oplus k} \oplus M'$. After a change of basis
in our presentation above, we may assume the first $n - k$
basis vectors of $R^{\oplus n}$ map into the summand $M'$ of $M$
and the last $k$-basis vectors of $R^{\oplus n}$ map to
basis elements of the summand $R^{\oplus k}$ of $M$.
Having done so, the last $k$ rows of the matrix $A$ vanish.
In this way we see that, replacing $M$ by $M'$, $k$ by $0$,
$n$ by $n - k$, and $A$ by the submatrix where we
delete the last $k$ rows, we reduce to the case discussed
in the next paragraph.

\medskip\noindent
Assume $R$ is local, $k = 0$, and $M$ annihilated by $f$.
Now the $0$th Fitting ideal of $M$ is $(f)$ and is generated by the
$n \times n$ minors of the matrix $A$ of size $n \times m$.
(This in particular implies $m \geq n$.)
Since $R$ is local, some $n \times n$ minor of $A$ is $uf$
for a unit $u \in R$.
After renumbering we may assume this minor is the first one.
Moreover, we know all other $n \times n$ minors of $A$ are
divisible by $f$. Write $A = (A_1 A_2)$ in block form where
$A_1$ is an $n \times n$ matrix and $A_2$ is an
$n \times (m - n)$ matrix. By
Algebra, Lemma \ref{algebra-lemma-matrix-right-inverse}
applied to the transpose of $A$ (!) we find there exists an
$n \times n$ matrix $B$ such that
$$
BA = B(A_1 A_2) = f
\left(
\begin{matrix}
u 1_{n \times n} & C
\end{matrix}
\right)
$$
for some $n \times (m - n)$ matrix $C$ with coefficients in $R$.
Then we first conclude $BA_1 = fu 1_{n \times n}$.
Thus
$$
BA_2 = fC = u^{-1}fuC = u^{-1}BA_1C
$$
Since the determinant of $B$ is a nonzerodivisor we conclude
that $A_2 = u^{-1}A_1C$. Therefore the image of $A$ is equal
to the image of $A_1$ which is isomorphic to $R^{\oplus n}$
because the determinant of $A_1$ is a nonzerodivisor.
Hence $M$ has projective dimension $\leq 1$.

\medskip\noindent
We return to the case of a general ring $R$. By the local case
we see that $M/M'$ is a finite locally free module of rank $k$, see
Algebra, Lemma \ref{algebra-lemma-finite-projective}. Hence the extension
$0 \to M' \to M \to M/M' \to 0$ splits. It follows that $M'$
is a finitely presented module. Choose a short exact sequence
$0 \to K \to R^{\oplus a} \to M' \to 0$. Then $K$ is a finite
$R$-module, see Algebra, Lemma \ref{algebra-lemma-extension}.
By the local case we see that $K_\mathfrak p \cong R_\mathfrak p^{\oplus a}$
for all primes. Hence by
Algebra, Lemma \ref{algebra-lemma-finite-projective}
again we see that $K$ is finite locally free of rank $a$.
It follows that $M'$ has projective dimension $\leq 1$
and the lemma is proved.
\end{proof}